\section{はじめに}
画像データセットに対する物体アノテーションは，画像中のどこに物体が存在するか（物体領域）と，その物体が何であるか（カテゴリ）という，2種類の異なる判断から構成される．
YOLO\cite{ref_YOLO}やDetection Transformer（DETR）\cite{ref_DETR}に代表される一般的な物体検出モデルは，これら2つの判断を単一のモデルに統合し，あらかじめ定めた固定のクラス集合に対して同時に学習する．
そのため，学習時に定めたクラス集合の範囲内では高精度な検出が可能である一方，新たなカテゴリを追加したり，カテゴリの粒度を変更したりする場合には，モデル自体の再学習が必要になる．
例えば，COCO\cite{ref_COCO}の80クラスで学習したYOLOは，ペンギンをbird，ラクダをdogやhorseのように80クラスのいずれかの語でしか表せず，80クラスのいずれにも似ていない物体はそもそも検出しない．

本研究では，物体領域の検出とカテゴリのアノテーションを明確に分離し，検出した領域のクロップに視覚言語モデル（Vision-Language Model，VLM）のゼロショット分類を適用する．
分離により，一度検出した領域に対し，候補語彙を変えるだけでカテゴリのアノテーションを繰り返すことができる．
この性質を用い，WordNet\cite{ref_WordNet}の上位語・下位語の関係をたどることで，追加学習なしに「dog → hunting dog → hound」のように粒度を変更できるアノテーションを実現する．

本稿では，この枠組みを3つの構成で評価する．
第一は，クラス非依存の検出にDETRを用い，深さ0のカテゴリも含めてSigLIP\cite{ref_SigLIP}で推定する提案手法である．
第二は，COCOの範囲で高精度なYOLOの予測を深さ0とし，YOLOのボックスのクロップに対してSigLIPとWordNetで粒度を変更するハイブリッド手法である．
第三は，ハイブリッド手法が検出しないCOCO外の物体を，物体らしさに基づくクラス非依存の検出器で補い，補った領域にも命名と粒度の変更を行う構成である．
比較対象として，YOLO単体を用いる．

本稿の貢献は以下の3点である．
第一に，COCO 2017 Validationと，WordNetのSynsetが付与されたLVIS\cite{ref_LVIS}のアノテーションを用い，YOLO単体，提案手法，ハイブリッド手法の深さ0の精度と，細かい粒度における階層的な精度を，検証用とテスト用に分けた画像で比較する．
第二に，COCO外の物体を見落とさないための検出器として，YOLO，DETR，OWLv2\cite{ref_OWLv2}の物体らしさのスコア，およびGrounding DINO\cite{ref_GroundingDINO}を，1画像あたりのボックス数をそろえて比較し，検出器がCOCO外の物体の粒度変更の上限を決めることを示す．
第三に，YOLOが検出しない領域のみを別の検出器で補う構成により，COCOの物体に対するYOLOの精度を保ったまま，COCO外の物体を検出から粒度の変更まで扱えることを示し，残る課題が命名の語彙にあることを明らかにする．
\add{さらに，加えた候補の棄却と，粗い語から順に選ぶ命名により，COCO外の物体の命名を改善できることを示す．}

本稿の構成は以下のとおりである．
2章で関連研究，3章で提案手法を述べ，4章で実験条件と結果を示す．
5章で考察を行い，6章で本稿をまとめる．
